\chapter{Kiến trúc hệ thống}
\label{chap:architecture}

Chương này trình bày kiến trúc hệ thống Self-Evolving AI Agents được xây dựng trên nền tảng NIKA. Mục tiêu của hệ thống là cho phép agent tương tác với môi trường mạng giả lập, thu thập bằng chứng chẩn đoán, đưa ra kết luận detection/localization/RCA và cải thiện hiệu quả qua các incident liên tiếp.

\section{Kiến trúc tổng quan}
\label{sec:overall_arch}

Hệ thống gồm bốn thành phần chính:
\begin{enumerate}
    \item \textbf{Môi trường mạng NIKA/Kathará:} Khởi tạo các topology mạng ảo hóa bằng Docker container, hỗ trợ nhiều giao thức và công nghệ như BGP, OSPF, RIP, SDN, P4/BMv2, DNS, DHCP và HTTP.
    \item \textbf{AI Agent chẩn đoán:} Sử dụng LLM làm lõi suy luận, vận hành theo workflow ReAct, Plan \& Execute hoặc Reflexion trên LangGraph \cite{langgraph2024}.
    \item \textbf{MCP Servers:} Cung cấp giao diện chuẩn hóa để agent gọi công cụ kiểm tra kết nối, cấu hình host, bảng định tuyến, trạng thái dịch vụ, log P4/BMv2, telemetry và nộp kết quả chẩn đoán \cite{mcp2024}.
    \item \textbf{Mô-đun tự tiến hóa:} Bao gồm Procedural Memory để lưu và truy xuất kỹ năng chẩn đoán, cùng Tool Refinement để tinh chỉnh tài liệu công cụ.
\end{enumerate}

Trong mỗi incident, NIKA khởi động lab tương ứng, tiêm lỗi vào môi trường mạng, sau đó agent thực hiện chẩn đoán bằng cách gọi các công cụ MCP. Kết quả cuối cùng được gửi qua Task Management MCP Server dưới dạng \texttt{is\_anomaly}, \texttt{faulty\_devices} và \texttt{root\_cause\_name}. Sau khi episode kết thúc, hệ thống đánh giá kết quả và kích hoạt các mô-đun học ngoại tuyến.

\section{Workflow chẩn đoán}
\label{sec:agent_workflows}

Workflow chính được sử dụng trong thực nghiệm là ReAct. Agent luân phiên giữa bước suy luận và bước gọi công cụ: từ triệu chứng ban đầu, agent đặt giả thuyết, chọn công cụ kiểm chứng, đọc kết quả và cập nhật hướng điều tra. Giới hạn \texttt{max\_steps} giúp tránh vòng lặp vô hạn và kiểm soát chi phí.

Hệ thống cũng hỗ trợ Plan \& Execute và Reflexion. Plan \& Execute tách quá trình lập kế hoạch khỏi thực thi, phù hợp với các sự cố cần chuỗi điều tra dài. Reflexion bổ sung bước tự đánh giá và thử lại khi kết quả chưa đạt yêu cầu. Trong phạm vi đánh giá chính, ReAct được chọn làm workflow so sánh vì đơn giản, ổn định và là baseline phổ biến cho agent sử dụng công cụ.

\section{Procedural Memory}
\label{sec:memory_module}

Procedural Memory lưu trữ kinh nghiệm chẩn đoán dưới dạng \texttt{ProceduralSkill}. Mỗi kỹ năng gồm điều kiện kích hoạt, chuỗi bước thực hiện, điều kiện kết thúc và các thuộc tính ngữ cảnh như scenario, protocol, service, symptom và tool. Trạng thái bộ nhớ được lưu trong \texttt{runtime/memory/\{bank\_id\}/skills.json}. Thiết kế này lấy cảm hứng từ hướng procedural memory của Skill-Pro.

Trước mỗi episode, hệ thống sinh \texttt{MemoryQuery} từ ngữ cảnh public của incident và truy xuất các kỹ năng phù hợp. Điểm truy xuất kết hợp độ khớp từ khóa, độ khớp scenario, độ khớp thuộc tính và độ tin cậy của kỹ năng. Trong quá trình agent gọi công cụ, \texttt{SkillToolRuntime} có thể chèn hướng dẫn ngắn từ kỹ năng đang hoạt động để agent chọn đúng công cụ, đúng tham số và biết khi nào cần dừng hoặc chuyển giả thuyết.

Sau episode, hệ thống tính phần thưởng dựa trên detection, localization, RCA và chi phí thao tác. Các episode tốt được đưa vào experience pool. Từ đó, LLM sinh semantic gradient để đề xuất cải thiện kỹ năng; ứng viên kỹ năng mới chỉ được chấp nhận khi vượt qua PPO-style gate \cite{ppo2017}. Cơ chế này giúp bộ nhớ học từ ca thành công nhưng hạn chế việc ghi nhận kinh nghiệm nhiễu hoặc quá khớp.

\section{Tool Refinement}
\label{sec:tool_evolution}

Tool Refinement cải thiện cách agent hiểu và sử dụng công cụ MCP bằng cách tinh chỉnh tài liệu mô tả công cụ, không thay đổi mã nguồn thực thi. Trạng thái được lưu tại \texttt{runtime/tool\_evolution/\{library\_id\}/state.json}. Thiết kế này lấy cảm hứng từ cơ chế tinh chỉnh tài liệu công cụ của DRAFT.

Sau mỗi session, hệ thống trích xuất các lần gọi công cụ từ \texttt{messages.jsonl}, xác định các khoảng trống hiểu biết như thiếu mô tả tham số, sai tên thiết bị, thiếu điều kiện tiền đề hoặc chưa biết cách diễn giải output. Analyzer sinh gợi ý cải thiện, Rewriter viết lại mô tả công cụ với ràng buộc tham số, ví dụ đúng/sai và ghi chú diễn giải kết quả. Khi tài liệu đạt ngưỡng hội tụ, công cụ được đóng băng để tránh viết lại không cần thiết.

Hai mô-đun Procedural Memory và Tool Refinement bổ trợ cho nhau: Procedural Memory cải thiện chiến lược chẩn đoán ở mức quy trình, còn Tool Refinement cải thiện độ chính xác khi sử dụng từng công cụ cụ thể.
